\documentclass[11pt]{article}
\usepackage[margin=1in]{geometry}
\usepackage{amsmath,amssymb}
\usepackage{listings}
\usepackage{xcolor}
\usepackage{hyperref}

\lstset{
    basicstyle=\ttfamily\footnotesize,
    breaklines=true,
    frame=single,
    language=Julia,
    keywordstyle=\color{blue},
    commentstyle=\color{gray},
    stringstyle=\color{teal},
    showstringspaces=false
}

\title{Boundary Conditions in JEXPRESSO\\
\large Code structure and user-side workflow}
\author{}
\date{}

\begin{document}
\maketitle

\section{Overview}

In JEXPRESSO the application of boundary conditions (BCs) is split between
two layers:

\begin{enumerate}
    \item A \textbf{kernel layer} located in
    \texttt{src/kernel/boundaryconditions/}, which contains the generic
    machinery that loops over boundary edges (2D) or faces (3D), reads the
    \texttt{tag} associated with every boundary entity in the Gmsh mesh,
    fetches normals, and pushes the boundary-node values either into
    \texttt{uaux} (for time-dependent problems with explicit time
    integration) or directly into the global linear system $L\,x = b$
    (for elliptic problems).
    \item A \textbf{user layer} located in
    \texttt{problems/<EQUATION>/<CASE>/user\_bc.jl}, where the user
    prescribes the actual boundary values or boundary fluxes by writing
    one (or two) Julia functions:
    \texttt{user\_bc\_dirichlet!} (essential conditions) and
    \texttt{user\_bc\_neumann}/\texttt{user\_bc\_neumann!} (natural
    conditions / surface fluxes).
\end{enumerate}

The bridge between the two layers is provided by
\texttt{src/kernel/boundaryconditions/custom\_bcs.jl}, which contains the
thin wrappers \texttt{dirichlet!} and \texttt{neumann} that simply
forward the call to the user-defined \texttt{user\_bc\_*} functions.

The high-level driver in \texttt{problems/drivers.jl} chooses
\emph{which} BC entry point is invoked according to the type of problem
(hyperbolic with explicit time integration, elliptic $A x = b$, or
hyperbolic with implicit/IMEX time integration that places terms on the
left-hand side).

\section{Kernel-side structure}

\subsection{Files in \texttt{src/kernel/boundaryconditions/}}

\begin{description}
    \item[\texttt{BCs.jl}] Defines the public entry points called by the
    driver and the time integrators:
    \begin{itemize}
        \item \texttt{apply\_boundary\_conditions\_dirichlet!}
        \item \texttt{apply\_boundary\_conditions\_neumann!}
        \item \texttt{apply\_boundary\_conditions\_lin\_solve!}
        \item \texttt{apply\_periodicity!}
    \end{itemize}
    Each public entry point dispatches to a dimension-specific
    \texttt{build\_custom\_bcs\_*!} method using the spatial-dimension
    type (\texttt{NSD\_1D}, \texttt{NSD\_2D}, \texttt{NSD\_3D}). For the
    linear-solve case there is in addition a sparse variant
    (\texttt{build\_custom\_bcs\_lin\_solve\_sparse!}) that mutates a
    \texttt{SparseMatrixCSC}.
    \item[\texttt{custom\_bcs.jl}] Wrappers \texttt{dirichlet!} and
    \texttt{neumann} that simply forward to the
    user-defined \texttt{user\_bc\_dirichlet!} /
    \texttt{user\_bc\_neumann}.
    \item[\texttt{surface\_integral.jl}] Helpers that build edge / face
    mass matrices and assemble surface integrals
    (\texttt{compute\_surface\_integral!},
    \texttt{compute\_segment\_integral!}, \texttt{DSS\_*\_integral!},
    \texttt{bulk\_surface\_flux!}). They are used by the Neumann path to
    add the boundary surface flux contribution to the RHS:
    \[
        \text{RHS}_i \mathrel{+}= \int_{\partial \Omega} F_{\text{surf}}\,\psi_i\,dS .
    \]
\end{description}

\subsection{Dirichlet machinery}

For 2D, \texttt{build\_custom\_bcs\_dirichlet!(::NSD\_2D, \ldots)} loops
over every boundary edge \texttt{iedge}; for each edge it skips
\texttt{periodicx}, \texttt{periodicz}, \texttt{periodic1},
\texttt{periodic2} and \texttt{Laguerre} tags, and for the remaining
$ngl$ Lagrange-Gauss-Lobatto (LGL) nodes it
\begin{enumerate}
    \item fills a sentinel value \texttt{qbdy[ieq]\,=\,4325789.0};
    \item calls \texttt{user\_bc\_dirichlet!(@view(uaux[ip,:]),\ldots,
    bdy\_edge\_type[iedge], qbdy, nx, ny, @view(qe[ip,:]),
    inputs[:SOL\_VARS\_TYPE])};
    \item if the user changed the sentinel and the value differs from the
    current solution, it overwrites \texttt{uaux[ip,ieq]} and zeroes
    \texttt{RHS[ip,ieq]} (so the time integrator does not modify the
    constrained DOF).
\end{enumerate}
At the end the array \texttt{u} is repacked from \texttt{uaux} via
\texttt{uaux2u!(u, uaux, neqs, npoin)}.

The 3D version is analogous, looping on \texttt{nfaces\_bdy} and using
\texttt{poin\_in\_bdy\_face} together with the 3-component normal
$(n_x,n_y,n_z)$.

\subsection{Neumann machinery}

\texttt{build\_custom\_bcs\_neumann!} is invoked only when
\texttt{inputs[:bdy\_fluxes]\,=\,true}. For each boundary edge/face it

\begin{enumerate}
    \item zeroes \texttt{F\_surf};
    \item if \texttt{inputs[:bulk\_fluxes]} is true, calls the analytic
    \texttt{bulk\_surface\_flux!}; for the special tag \texttt{"MOST"}
    in 3D it calls the Monin-Obukhov routine \texttt{CM\_MOST!};
    \item otherwise calls the user routine
    \texttt{user\_bc\_neumann!(@view(F\_surf[\ldots]), uaux[ip,:],
    uaux[ip1,:], qe[ip,:], qe[ip1,:], bdy\_edge\_type[iedge],
    coords, \ldots, inputs[:SOL\_VARS\_TYPE])};
    \item integrates over the boundary segment/face with
    \texttt{compute\_segment\_integral!} or
    \texttt{compute\_surface\_integral!};
    \item performs DSS and adds the result to the global RHS,
    \texttt{RHS[:,ieq] .+= S\_flux[:,ieq]}.
\end{enumerate}

\subsection{Linear-solve machinery}

The variant used by elliptic problems is
\texttt{apply\_boundary\_conditions\_lin\_solve!}, which dispatches to
either the dense or the sparse kernel:

\begin{itemize}
    \item \texttt{build\_custom\_bcs\_lin\_solve!(::NSD\_2D, \ldots)} (dense
    \texttt{Matrix\{Float64\}}). For every Dirichlet boundary node
    \texttt{ip} it
    \begin{enumerate}
        \item calls \texttt{user\_bc\_dirichlet!} \emph{but writes the
        prescribed value into \texttt{RHS[ip,:]}} (and into the helper
        \texttt{qbdy});
        \item zeroes the row \texttt{L[ip, :]\,=\,0};
        \item sets \texttt{L[ip, ip]\,=\,1};
        \item assigns \texttt{RHS[ip, ieq]\,=\,qbdy[ieq]}.
    \end{enumerate}
    This is the standard ``identity row'' enforcement of essential BCs in
    $A x = b$.
    \item \texttt{build\_custom\_bcs\_lin\_solve\_sparse!} does the same on
    a \texttt{SparseMatrixCSC}, calling
    \texttt{apply\_dirichlet\_bc\_inplace!} which collects the unique
    boundary nodes from \texttt{poin\_in\_bdy\_edge} and rewrites
    \texttt{L[ip,:]\,=\,0; L[ip,ip]\,=\,1}.
\end{itemize}

\section{Driver-side dispatch}

The function \texttt{driver(\ldots)} in \texttt{problems/drivers.jl}
performs the high-level dispatch:

\begin{lstlisting}
if !inputs[:llinsolve]
    # Hyperbolic / parabolic:  M dq/dt = RHS
    solution = time_loop!(inputs, params, u, partitioned_model)
else
    # Elliptic:  L*q = M*RHS
    ...
    apply_boundary_conditions_lin_solve!(sem.matrix.L, ..., RHS, ...)
    sol = sem.matrix.L \ RHS
    write_output(...)
end
\end{lstlisting}

Therefore the same \texttt{user\_bc\_dirichlet!} that a user writes in
\texttt{user\_bc.jl} is reused in three different contexts:

\begin{enumerate}
    \item \emph{Hyperbolic with explicit time integration} -- called from
    inside \texttt{time\_loop!} via
    \texttt{apply\_boundary\_conditions\_dirichlet!} (and
    \texttt{\dots\_neumann!} when boundary fluxes are activated). The
    constraint is enforced on \texttt{uaux} and the corresponding row of
    the RHS is zeroed so that the explicit RK step does not drift the
    boundary values.
    \item \emph{Elliptic ($A x = b$)} -- called from
    \texttt{apply\_boundary\_conditions\_lin\_solve!}. The same user
    routine is used, but the kernel writes the prescribed value into
    \texttt{RHS} and modifies the system matrix \texttt{L} (zeroing the
    row and putting a 1 on the diagonal).
    \item \emph{Hyperbolic with implicit / IMEX time integration} -- the
    implicit operator is assembled in the same Laplace-style matrix
    \texttt{sem.matrix.L} that the elliptic path uses; the same
    \texttt{apply\_boundary\_conditions\_lin\_solve!} entry point is
    invoked at every implicit stage to constrain the rows that
    correspond to Dirichlet nodes. Neumann contributions are still
    assembled by \texttt{apply\_boundary\_conditions\_neumann!} and
    added to the right-hand side.
\end{enumerate}

\section{User-side workflow}

A user creates a problem under
\texttt{problems/<EQUATION>/<CASE>/} and is required to provide
\texttt{user\_bc.jl}. The Gmsh \texttt{.msh} file declares physical
groups whose names become the \texttt{tag} string passed to the BC
function. Inside \texttt{user\_bc.jl} the user uses an \texttt{if/elseif}
ladder on \texttt{tag} to set the prescribed values:

\begin{lstlisting}
if (tag == "inflow")
    qbdy[1] = 3.0
elseif (tag == "fix_temperature")
    qbdy[2] = 300.0
end
\end{lstlisting}

Any component that the user does not assign keeps the sentinel
\texttt{4325789.0}, which the kernel detects and \emph{leaves
unchanged}, i.e.\ this is interpreted as a ``do nothing'' (natural)
condition for that component.

\subsection{Hyperbolic problems with explicit time integration}

Example: \texttt{problems/CompEuler/theta/} (rising thermal bubble).

\paragraph{user\_inputs.jl.} The user selects an explicit ODE solver
and leaves \texttt{:llinsolve} unset (i.e.\ false, the default), which
selects the \texttt{time\_loop!} branch in the driver:
\begin{lstlisting}
:ode_solver  => CarpenterKennedy2N54(),
:Delta t     => 0.5,
:tinit       => 0.0,
:tend        => 1000.0,
:SOL_VARS_TYPE => PERT(),
\end{lstlisting}

\paragraph{user\_bc.jl.} A free-slip wall is implemented by
projecting out the normal component of the momentum
$(\rho u, \rho v)$. Both the \texttt{TOTAL} and the \texttt{PERT}
variable formulations are provided:

\begin{lstlisting}
function user_bc_dirichlet!(q, coords, t, tag, qbdy, nx, ny, qe, ::PERT)
    qnl = nx*(q[2]+qe[2]) + ny*(q[3]+qe[3])
    qbdy[2] = (q[2]+qe[2] - qnl*nx) - qe[2]
    qbdy[3] = (q[3]+qe[3] - qnl*ny) - qe[3]
end
\end{lstlisting}

The Neumann routine returns a zero flux because no surface flux is
needed for this case:
\begin{lstlisting}
function user_bc_neumann(q, gradq, coords, t, tag, inputs::Dict)
    return zeros(size(q,2),1)
end
\end{lstlisting}

\paragraph{What the kernel does at every RK stage.} The integrator
calls \texttt{apply\_boundary\_conditions\_dirichlet!}, which for every
boundary node \texttt{ip} invokes the user routine with the local
normal $(n_x,n_y)$ taken from \texttt{metrics.nx,~metrics.ny}, copies
the returned \texttt{qbdy} into \texttt{uaux[ip,:]} and zeros
\texttt{RHS[ip,:]}. The result is then mapped back into \texttt{u} via
\texttt{uaux2u!}.

\subsection{Elliptic problems leading to $A x = b$}

Example: \texttt{problems/Elliptic/case1/} (Laplace with non-homogeneous
Dirichlet data on a $5\times\pi$ rectangle).

\paragraph{user\_inputs.jl.} The user activates the linear-solve path:
\begin{lstlisting}
:llinsolve            => true,
:ldss_laplace         => true,
:rconst               => [0.0],
\end{lstlisting}
which forces the driver into the \texttt{else} branch where the
operator \texttt{sem.matrix.L} (the Laplacian) is assembled and the
system \texttt{L * q = M * RHS} is solved by the chosen Krylov method
(\texttt{:ode\_solver => "BICGSTABLE"}) or by the direct backslash.

\paragraph{user\_bc.jl.} Only \texttt{user\_bc\_dirichlet!} is needed:
\begin{lstlisting}
function user_bc_dirichlet!(q, coords, t, tag, qbdy, nx, ny, qe, ::TOTAL)
    L = 5.0
    if     (tag == "bottom") qbdy[1] = 100.0
    elseif (tag == "right")  qbdy[1] = 100.0
    elseif (tag == "top")    qbdy[1] = 100 + 20.0*sin(pi*coords[1]/L)
    elseif (tag == "left")   qbdy[1] = 100.0
    end
end
\end{lstlisting}

\paragraph{What the kernel does once.} Before the linear solve, the
driver calls
\texttt{apply\_boundary\_conditions\_lin\_solve!(sem.matrix.L, \ldots,
RHS, \ldots)}. For every node on a non-periodic boundary, the user
routine fills \texttt{qbdy}; the kernel then sets

\[
   L[ip,:] = 0, \qquad L[ip,ip] = 1, \qquad
   \text{RHS}[ip] = \texttt{qbdy[1]} ,
\]

so that the linear system enforces $q_{ip} = \texttt{qbdy[1]}$
exactly. Internally this is exactly the loop in
\texttt{build\_custom\_bcs\_lin\_solve!} (or its sparse counterpart
when \texttt{:lsparse => true}).

\subsection{Hyperbolic problems with implicit / IMEX time integration}

When part of the operator (typically a stiff diffusion or acoustic
operator) is treated implicitly, that operator is assembled into the
same matrix \texttt{sem.matrix.L} that the elliptic path uses, and the
implicit (or IMEX) stage solves
\[
   (M + \Delta t\,\theta\,L)\,q^{\,n+1} \;=\; M\,q^{\,n} + \Delta t\,R(q^{\,n}, q^{\,n+1}).
\]
At every implicit stage the BC kernel is therefore called twice:

\begin{enumerate}
    \item \texttt{apply\_boundary\_conditions\_lin\_solve!} on the
    implicit operator and on the right-hand side. This is the same
    routine described in section 4.2 above: rows of the implicit
    operator corresponding to Dirichlet nodes are replaced by an
    identity row and the prescribed value is written into the RHS, so
    that the Krylov solver returns the exact boundary value at each
    implicit stage.
    \item \texttt{apply\_boundary\_conditions\_neumann!}, which, when
    \texttt{inputs[:bdy\_fluxes] = true}, evaluates the user-defined
    surface flux through \texttt{user\_bc\_neumann!}, performs the
    surface quadrature and the DSS, and adds the result to the
    explicit part of the RHS via
    \[
        \text{RHS}[:,ieq] \mathrel{+}= S_{\text{flux}}[:,ieq].
    \]
\end{enumerate}

\paragraph{User responsibility.} The user does \emph{not} have to
rewrite the BC code: the same \texttt{user\_bc\_dirichlet!} and
\texttt{user\_bc\_neumann!} written for the explicit case are reused
verbatim. The user only has to:
\begin{itemize}
    \item select an implicit or IMEX integrator in
    \texttt{user\_inputs.jl};
    \item set \texttt{:llinsolve\,=>\,true} when the chosen integrator
    requires the assembly of the implicit operator into
    \texttt{sem.matrix.L};
    \item set \texttt{:bdy\_fluxes\,=>\,true} (and optionally
    \texttt{:bulk\_fluxes\,=>\,true} for the built-in bulk surface
    flux, or \texttt{tag\,=\,"MOST"} for Monin-Obukhov) if surface
    fluxes have to be added to the RHS at every stage.
\end{itemize}

\section{Summary table}

\begin{center}
\begin{tabular}{|l|l|l|l|}
\hline
Problem class & Driver branch & Kernel entry point & Action \\
\hline
Hyperbolic, explicit  & \texttt{time\_loop!} & \texttt{apply\_boundary\_conditions\_dirichlet!} & overwrite \texttt{uaux}, zero RHS row \\
                      &                      & \texttt{apply\_boundary\_conditions\_neumann!}   & add surface flux to RHS \\
\hline
Elliptic ($Ax=b$)     & \texttt{llinsolve} branch & \texttt{apply\_boundary\_conditions\_lin\_solve!} & zero $L$ row, $L_{ii}=1$, write RHS \\
\hline
Hyperbolic, IMEX/impl.& \texttt{time\_loop!} + implicit stage & \texttt{apply\_boundary\_conditions\_lin\_solve!} & same as elliptic on stage matrix \\
                      &                      & \texttt{apply\_boundary\_conditions\_neumann!}   & add surface flux to RHS \\
\hline
\end{tabular}
\end{center}

\section{Files referenced}

\begin{itemize}
    \item \texttt{src/kernel/boundaryconditions/BCs.jl}
    \item \texttt{src/kernel/boundaryconditions/custom\_bcs.jl}
    \item \texttt{src/kernel/boundaryconditions/surface\_integral.jl}
    \item \texttt{problems/drivers.jl}
    \item \texttt{problems/CompEuler/theta/user\_bc.jl},
          \texttt{user\_inputs.jl}
    \item \texttt{problems/Elliptic/case1/user\_bc.jl},
          \texttt{user\_inputs.jl}
\end{itemize}

\end{document}
